% GENERATED FILE. Edit canonical lessons, then run npm run generate:books.
% canonical-source-hash: fnv1a64:65aaa9ed6f083c31
% canonical-lessons: LA-C233-insto, LA-C233-interrumpo, LA-C233-ligo, LA-C233-misceo, LA-C233-pareo

\chapter{To press on, To break off, To tie, To mix, To obey}
\label{ch:la-to-press-on-to-break-off-to-tie-to-mix-to-obey}
\hlchaptermodality{233}

\begin{chapteropening}
\textbf{By the end of this chapter:} \emph{I can say to press on, to break off, to tie, to mix and to obey.}

Complete the last lesson of chapter 233: I can say to press on, to break off, to tie, to mix and to obey.
\end{chapteropening}

\section[īnstō, īnstāre]{īnstō, īnstāre — to press on, to insist}

\label{lesson:LA-C233-insto}



\begin{quote}
Before the new one: say the Latin for to serve, then the Latin for to suffer.
\end{quote}

\subsection*{You'll want to know: īnstō, īnstāre}
\textbf{īnstō, īnstāre} — \textquotedblleft{}to press on, to insist\textquotedblright{}.

\begin{cousinweb}
\textbf{īnstō} is \emph{in}, \textquotedblleft{}on, upon\textquotedblright{}, with \emph{stō}, \textquotedblleft{}to stand\textquotedblright{} (EtymDB 2.1): to stand upon, to press on. English \emph{instant} comes from its participle.
\end{cousinweb}

\begin{grammarlens}[title={What you've built}]
One more word for this chapter.
\end{grammarlens}

\subsection*{Your turn}
\begin{itemize}
\raggedright
  \item \emph{Say it:} \emph{īnstō, īnstāre}
  \item \emph{Say it:} \emph{īnstō, īnstāre}, once more
  \item \emph{Say it:} say \emph{īnstō, īnstāre}
  \item \emph{Recall:} say the Latin for to serve, then the Latin for to suffer, then say \emph{īnstō, īnstāre} again
\end{itemize}

\begin{tcolorbox}[breakable,colback=teal!4,colframe=teal!35!black,title={Before you move on}]
What does īnstō, īnstāre mean? (\textquotedblleft{}To press on, to insist\textquotedblright{}.) Say it once more.
\end{tcolorbox}
\section[interrumpō, interrumpere]{interrumpō, interrumpere — to break off, to interrupt}

\label{lesson:LA-C233-interrumpo}



\begin{quote}
Before the new one: say the Latin for to suffer, then the Latin for to press on.
\end{quote}

\subsection*{You'll want to know: interrumpō, interrumpere}
\textbf{interrumpō, interrumpere} — \textquotedblleft{}to break off, to interrupt\textquotedblright{}.

\begin{cousinweb}
\textbf{interrumpō} is \emph{inter}, \textquotedblleft{}between\textquotedblright{}, with \emph{rumpō}, \textquotedblleft{}to break\textquotedblright{} (EtymDB 2.1): to break in between. \emph{rumpō} is usually traced to Proto-Indo-European \emph{*Hrewp-}, \textquotedblleft{}to break\textquotedblright{}.
\end{cousinweb}

\begin{grammarlens}[title={What you've built}]
One more word for this chapter.
\end{grammarlens}

\subsection*{Your turn}
\begin{itemize}
\raggedright
  \item \emph{Say it:} \emph{interrumpō, interrumpere}
  \item \emph{Say it:} \emph{interrumpō, interrumpere}, once more
  \item \emph{Say it:} say \emph{interrumpō, interrumpere}
  \item \emph{Recall:} say the Latin for to suffer, then the Latin for to press on, then say \emph{interrumpō, interrumpere} again
\end{itemize}

\begin{tcolorbox}[breakable,colback=teal!4,colframe=teal!35!black,title={Before you move on}]
What does interrumpō, interrumpere mean? (\textquotedblleft{}To break off, to interrupt\textquotedblright{}.) Say it once more.
\end{tcolorbox}
\section[ligō, ligāre]{ligō, ligāre — to tie}

\label{lesson:LA-C233-ligo}



\begin{quote}
Before the new one: say the Latin for to press on, then the Latin for to break off.
\end{quote}

\subsection*{You'll want to know: ligō, ligāre}
\textbf{ligō, ligāre} — \textquotedblleft{}to tie\textquotedblright{}.

\begin{cousinweb}
\textbf{ligō} is probably inherited, from a Proto-Indo-European root \emph{*leyg-}, \textquotedblleft{}to bind\textquotedblright{} (EtymDB 2.1). English \emph{ligament} comes from it.
\end{cousinweb}

\begin{grammarlens}[title={What you've built}]
One more word for this chapter.
\end{grammarlens}

\subsection*{Your turn}
\begin{itemize}
\raggedright
  \item \emph{Say it:} \emph{ligō, ligāre}
  \item \emph{Say it:} \emph{ligō, ligāre}, once more
  \item \emph{Say it:} say \emph{ligō, ligāre}
  \item \emph{Recall:} say the Latin for to press on, then the Latin for to break off, then say \emph{ligō, ligāre} again
\end{itemize}

\begin{tcolorbox}[breakable,colback=teal!4,colframe=teal!35!black,title={Before you move on}]
What does ligō, ligāre mean? (\textquotedblleft{}To tie\textquotedblright{}.) Say it once more.
\end{tcolorbox}
\section[misceō, miscēre]{misceō, miscēre — to mix}

\label{lesson:LA-C233-misceo}



\begin{quote}
Before the new one: say the Latin for to break off, then the Latin for to tie.
\end{quote}

\subsection*{You'll want to know: misceō, miscēre}
\textbf{misceō, miscēre} — \textquotedblleft{}to mix\textquotedblright{}.

\begin{cousinweb}
\textbf{misceō} is inherited, usually traced to Proto-Indo-European \emph{*meyk-}, \textquotedblleft{}to mix\textquotedblright{} (EtymDB 2.1); Greek \emph{mignymi}, \textquotedblleft{}to mix\textquotedblright{}, is a cousin (Lewis \& Short). English \emph{mix} comes from its participle \emph{mixtus}, \textquotedblleft{}mixed\textquotedblright{}.
\end{cousinweb}

\begin{grammarlens}[title={What you've built}]
One more word for this chapter.
\end{grammarlens}

\subsection*{Your turn}
\begin{itemize}
\raggedright
  \item \emph{Say it:} \emph{misceō, miscēre}
  \item \emph{Say it:} \emph{misceō, miscēre}, once more
  \item \emph{Say it:} say \emph{misceō, miscēre}
  \item \emph{Recall:} say the Latin for to break off, then the Latin for to tie, then say \emph{misceō, miscēre} again
\end{itemize}

\begin{tcolorbox}[breakable,colback=teal!4,colframe=teal!35!black,title={Before you move on}]
What does misceō, miscēre mean? (\textquotedblleft{}To mix\textquotedblright{}.) Say it once more.
\end{tcolorbox}
\section[pāreō, pārēre]{pāreō, pārēre — to obey}

\label{lesson:LA-C233-pareo}



\begin{quote}
Before the new one: say the Latin for to tie, then the Latin for to mix.
\end{quote}

\subsection*{You'll want to know: pāreō, pārēre}
\textbf{pāreō, pārēre} — \textquotedblleft{}to obey\textquotedblright{}.

\begin{cousinweb}
\textbf{pāreō} is probably inherited, from a Proto-Indo-European \emph{*peh2s-}, \textquotedblleft{}to watch, to see\textquotedblright{} (EtymDB 2.1): first to appear, to be at hand, then to be at someone's call, to obey.
\end{cousinweb}

\begin{grammarlens}[title={What you've built}]
One more word for this chapter.
\end{grammarlens}

\subsection*{Your turn}
\begin{itemize}
\raggedright
  \item \emph{Say it:} \emph{pāreō, pārēre}
  \item \emph{Say it:} \emph{pāreō, pārēre}, once more
  \item \emph{Say it:} say \emph{pāreō, pārēre}
  \item \emph{Recall:} say the Latin for to tie, then the Latin for to mix, then say \emph{pāreō, pārēre} again
\end{itemize}

\begin{tcolorbox}[breakable,colback=teal!4,colframe=teal!35!black,title={Before you move on}]
What does pāreō, pārēre mean? (\textquotedblleft{}To obey\textquotedblright{}.) Say it once more.
\end{tcolorbox}
